\chapter{\IfLanguageName{dutch}{Analyse en evaluatie}{Analysis and evaluation}}%
\label{ch:analyse-evaluatie}

Dit hoofdstuk analyseert de resultaten van de testfase en beantwoordt de
onderzoeksvragen zoals geformuleerd in de inleiding. De bevindingen worden
getoetst aan de literatuur en de doelstellingen van de proof of concept.
Aansluitend worden de beperkingen van het onderzoek geïdentificeerd en
aanbevelingen geformuleerd voor verdere doorontwikkeling.

\section{Beantwoording van de onderzoeksvragen}

\subsection{Welke moeilijkheden ervaren bezoekers bij het navigeren op kerkhoven?}

De literatuurstudie bevestigt dat bezoekers regelmatig moeilijkheden ervaren bij
het terugvinden van een specifiek graf, vooral op grote of uniform ingedeelde
begraafplaatsen \autocite{Haekkilae2022}. Het kerkhof van Sint-Gillis-Waas, met
zijn reguliere rasterindeling van identieke rijen en kolommen, illustreert dit
concreet: zonder digitale ondersteuning is het moeilijk om de juiste rij te
identificeren. De terreinanalyse bevestigde dat alle rijen visueel sterk op
elkaar lijken, wat een directe aanleiding was voor het ontwerp van de
verticale AR-grafmarker.

\subsection{Welke beperkingen hebben bestaande navigatietoepassingen?}

Bestaande oplossingen zoals het geoloket van de gemeente Beveren bieden
kaartnavigatie maar vereisen dat de gebruiker zijn blik afwisselt tussen het
scherm en de fysieke omgeving. Dit verhoogt de cognitieve belasting en past
minder goed bij de emotionele context van een begraafplaatsbezoek
\autocite{Zhao2025}. Het interne systeem WebBGP van Sint-Gillis-Waas is
bovendien uitsluitend beschikbaar voor gemeentepersoneel en biedt geen
directe ondersteuning aan bezoekers.

\subsection{Welke behoeften hebben bezoekers van verschillende leeftijden?}

De testresultaten tonen een duidelijk verschil in gebruikspatroon tussen de
twee externe testpersonen. In Scenario A slaagden beide testpersonen er niet
in het graf te vinden, ongeacht leeftijd of smartphoneervaring. Dit wijst erop
dat het probleem van kerkhofnavigatie zonder digitale hulp leeftijdsonafhankelijk
is. In Scenario C toonden de testpersonen echter een duidelijk verschil: TP2
(20 jaar, smartphoneervaring 4/5) navigeerde snel en intuïtief en ontdekte de
AR-functionaliteit zelfstandig, terwijl TP3 (49 jaar, smartphoneervaring 2/5)
een korte tip nodig had om de AR-overlay te ontdekken.

Dit bevestigt de bevindingen uit de literatuurstudie dat minder technologisch
vaardige gebruikers meer begeleiding nodig hebben bij het gebruik van AR-interfaces
\autocite{Haekkilae2022}. Beide testpersonen bereikten uiteindelijk het geselecteerde
graf en gaven een positief eindoordeel (8/10 en 7/10), wat aantoont dat de
basisfunctionaliteit voor beide leeftijdsgroepen toegankelijk is, mits een
minimale onboarding aanwezig is.

\subsection{Welke emotionele en contextuele factoren spelen een rol?}

Een begraafplaats vraagt om discretie en respect. De applicatie werd ontworpen
zonder geluid, zonder opvallende visuele effecten buiten het scherm en met een
interface die snel en efficiënt te bedienen is. De AR-overlay is uitsluitend
zichtbaar op het scherm van de gebruiker en heeft geen impact op de beleving
van andere bezoekers. Dit sluit aan bij de ethische overwegingen beschreven
in hoofdstuk~\ref{ch:juridische-analyse}.

\subsection{Welke privacygevoelige data is hierbij betrokken?}

De juridische analyse toonde aan dat namen, geboorte- en overlijdensdata van
overledenen buiten de beschermingssfeer van de GDPR vallen, maar dat zorgvuldig
gebruik vereist is. De locatiegegevens van de gebruiker zijn wel persoonsgegevens
en mogen niet worden opgeslagen. Het rijksregisternummer is categorisch
uitgesloten op basis van de Wet van 8 augustus 1983 \autocite{WetRijksregister1983}.

\subsection{Hoe kan Augmented Reality ondersteuning bieden bij kerkhofnavigatie?}

De vergelijking van de drie testscenario's toont de meerwaarde van AR het
duidelijkst aan. In Scenario A slaagden beide testpersonen er niet in het
graf te vinden. In Scenario B vonden ze het graf wel, maar misten ze
routebegeleiding. In Scenario C was de navigatie het meest volledig en
intuïtief: de 3D-navigatiepijl wees gedurende alle tests correct in de
richting van het geselecteerde graf en de grafmarker werd zichtbaar op circa
20 meter. AR biedt hierbij een fundamenteel voordeel ten opzichte van
klassieke kaartnavigatie: de gebruiker hoeft zijn blik niet weg te wenden
van de omgeving \autocite{Zhao2025}.

Opvallend is dat de AR-functionaliteit niet onmiddellijk voor alle
gebruikers duidelijk was. TP3 ontdekte de AR-overlay pas na een korte tip.
Dit wijst op de nood aan een minimale onboarding die nieuwe gebruikers
introduceert tot de werking van de applicatie.

\subsection{Hoe kan AR gecombineerd worden met digitale kaarten of geoloketten?}

De applicatie combineert beide benaderingen: de AR-cameraweergave voor directe
ruimtelijke navigatie en de satellietminikaart voor contextueel overzicht. De
testresultaten tonen aan dat beide elementen actief gebruikt werden: TP2 gebruikte
de minikaart als primaire oriëntatie, terwijl TP3 na het ontdekken van de AR
de kaart volledig liet vallen. De combinatie biedt dus flexibiliteit afhankelijk
van de voorkeur en het technologische profiel van de gebruiker.

In een productieomgeving kan de kaartcomponent worden uitgebreid met een geroute
navigatielijn op basis van OSM-paddata, zoals beschreven in
sectie~\ref{sec:routelijn}. Beide testpersonen gaven dit expliciet aan als
gewenste verbetering.

\subsection{Welke ontwerprichtlijnen zijn nodig om rust en sereniteit te bewaren?}

Op basis van de implementatie en de ethische analyse werden de volgende
ontwerprichtlijnen gehanteerd: geen geluidseffecten, een interface die snel
en discreet te bedienen is, AR-elementen die uitsluitend zichtbaar zijn op
het eigen scherm, en een eenvoudige gebruikersstroom zonder onnodige stappen.
Dit sluit aan bij de bevindingen van \autocite{Haekkilae2022} over de nood
aan respectvolle digitale oplossingen in emotioneel gevoelige contexten.

\subsection{Welke data kan veilig gebruikt worden?}

Naam, geboorte- en overlijdensdatum, GPS-locatie van het graf en
sectie/rij/plotnummer kunnen veilig worden verwerkt voor navigatiedoeleinden.
Het rijksregisternummer is absoluut uitgesloten. Biografieën en foto's vereisen
expliciete toestemming van nabestaanden. De GPS-locatie van de gebruiker mag
niet worden opgeslagen. Tabel~\ref{tab:datatypes} in
hoofdstuk~\ref{ch:juridische-analyse} geeft een volledig overzicht.

\subsection{Hoe kan deze proof of concept getest en geëvalueerd worden?}

De technische evaluatie toonde aan dat een combinatie van functionele controles
(pijlrichting, markerplaatsing, tracking, zoekfunctie) en observationele meting
(navigatietijd, gebruiksgemak, subjectief oordeel) een adequate evaluatiemethode
vormt voor een proof of concept in dit stadium. De toevoeging van externe
testpersonen met uiteenlopende technologische achtergronden leverde kwalitatieve
inzichten op die louter technische evaluatie niet had kunnen opleveren, zoals de
nood aan onboarding en de voorkeur voor de minikaart als primaire oriëntatie.

\section{Evaluatie van het proof of concept}

\subsection{Wat werkt goed}

De technische evaluatie en de gebruikerstests tonen aan dat de kernfunctionaliteit
correct werkt. De AR-navigatiepijl is betrouwbaar, responsief en wees gedurende
alle tests in de juiste richting. De grafmarker is zichtbaar van op een afstand
die praktisch bruikbaar is op kerkhofschaal. De zoekfunctie, minikaart en
afstandslabel functioneren allen correct. De applicatie opent onmiddellijk op
de cameraweergave zonder laadtijd of onboarding. Het fallback-mechanisme werkte
correct bij gesimuleerde tracking-degradatie. De subjectieve scores van 7/10 en
8/10 bevestigen dat de applicatie als bruikbaar en waardevol wordt ervaren.

\subsection{Wat kan beter}

Op basis van de testresultaten werden drie concrete verbeterpunten geïdentificeerd:

\begin{itemize}
    \item \textbf{Vorm van de navigatiepijl:} de huidige 3D-geometrie wordt
    niet onmiddellijk herkend als een navigatiepijl. Een meer herkenbare
    pijlvorm zou de bruikbaarheid verhogen, vooral voor gebruikers zonder
    AR-ervaring.
    
    \item \textbf{Ontbrekende onboarding:} geen van de externe testpersonen
    had eerder AR gebruikt, en TP3 ontdekte de AR-overlay pas na een tip.
    Een korte introductiemelding bij het eerste gebruik, die de gebruiker
    instrueert de camera rechtvooruit te richten, zou dit probleem oplossen.
    
    \item \textbf{Padsgewijze navigatie:} de pijl wijst steeds in een rechte
    lijn naar het geselecteerde graf. Beide externe testpersonen gaven aan dat
    een routelijn langs de bestaande paden, zowel op de kaart als in AR op de
    grond, de navigatie aanzienlijk zou verbeteren. Dit vereist GIS-paddata
    die momenteel niet beschikbaar is.
\end{itemize}

\section{Haalbaarheid als productieoplossing}

Op basis van de technische evaluatie en de juridische analyse kan worden
geconcludeerd dat het concept haalbaar is als productieoplossing voor de gemeente
Sint-Gillis-Waas, mits aan een aantal voorwaarden wordt voldaan:

\begin{enumerate}
    \item \textbf{Officiële grafdata:} de mockdataset moet worden vervangen door
    de reële grafdata van WebBGP via een API-koppeling onder een formele
    verwerkersovereenkomst.
    \item \textbf{Juridische afspraken:} een privacyverklaring, verwerkersovereenkomst
    en mogelijk een DPIA zijn vereiste stappen voorafgaand aan ingebruikname.
    \item \textbf{Platformuitbreiding:} een Android-versie via ARCore of een
    cross-platform implementatie via Unity AR Foundation vergroot het bereik.
    \item \textbf{Gebruikersstudie:} een formele test met diverse doelgroepen,
    inclusief oudere bezoekers, is noodzakelijk om effectiviteit en
    gebruiksvriendelijkheid volledig te bevestigen.
\end{enumerate}

\section{Aanbevelingen voor vervolgonderzoek}

\begin{itemize}
    \item \textbf{Onboarding toevoegen:} een korte introductiemelding bij het
    eerste gebruik die de gebruiker instrueert de camera rechtvooruit te richten.
    Dit was de meest directe en concrete aanbeveling van beide externe testpersonen.
    
    \item \textbf{Pijlvorm verbeteren:} de 3D-geometrie van de navigatiepijl
    aanpassen zodat ze beter overeenkomt met een klassieke, herkenbare pijlvorm.
    
    \item \textbf{Gebruikersstudie met diverse doelgroepen:} uitvoering van de
    drie vergelijkende testscenario's met testpersonen van verschillende leeftijden
    en technologische achtergrond, inclusief 65-plussers, om de effectiviteit van
    AR ten opzichte van bestaande navigatiemethoden kwantitatief te meten.
    
    \item \textbf{OSM-paddata integratie:} samenwerking met de lokale
    OpenStreetMap-gemeenschap om het kerkhof in kaart te brengen, zodat
    een geroute navigatielijn kan worden geïmplementeerd via een open-source
    routeringsbibliotheek \autocite{Novack2018}.
    
    \item \textbf{Koppeling met WebBGP:} onderzoek naar de technische en
    juridische haalbaarheid van een API-koppeling met het WebBGP-systeem
    van Cevi voor reële en officieel ingemeten grafdata.
    
    \item \textbf{Android-versie:} implementatie via ARCore of Unity AR
    Foundation om het bereik te vergroten naar Android-gebruikers
    \autocite{GoogleARCore2024}.
    
    \item \textbf{Evaluatie onder diverse weersomstandigheden:} testen bij
    bewolking, regen en avondlicht om de robuustheid van de ARKit tracking
    in minder gunstige omstandigheden te beoordelen.
\end{itemize}